% Part: first-order-logic
% Chapter: models-theories
% Section: expressing-relations

\documentclass[../../../include/open-logic-section]{subfiles}

\begin{document}

\olfileid{fol}{mat}{exr}

\olsection{\article{structure}
  \printtoken{S}{structure} मधील संबंध व्यक्त करणे}

\begin{explain}
!!{formula}s चा एक प्रमुख उपयोग म्हणजे !!a{structure}~$\Struct M$ मधील
गुणधर्म आणि संबंध हे~$\Struct M$ ची भाषा~$\Lang L$ हिच्या आदिम संकल्पनांच्या
साहाय्याने व्यक्त करणे. याचा अर्थ पुढीलप्रमाणे आहे: $\Struct M$ चा
!!{domain} हा वस्तूंचा संच असतो. !!{constant}s, !!{function}s आणि
!!{predicate}s यांचे~$\Struct M$ मध्ये अनुक्रमे~$\Domain M$ मधील काही वस्तू,
$\Domain M$ वरील फलने आणि~$\Domain M$ वरील संबंध यांनी अर्थनिर्धारण केलेले
असते. उदाहरणार्थ, $\Obj A^2_0$ हे~$\Lang L$ मध्ये असेल, तर~$\Struct M$
त्याला संबंध~$R = \Assign{{\Obj A^2_0}}{M}$ नेमते. मग सूत्र
$\Atom{\Obj A^2_0}{\Obj v_1, \Obj v_2}$ पुढील अर्थाने नेमका तो संबंध
\emph{व्यक्त करते}: चर-मूल्यांकन~$s$ हे~$\Obj v_1$ ला
$a \in \Domain{M}$ आणि~$\Obj v_2$ ला $b \in \Domain M$ शी संगत करत असेल,
तर
\[
Rab \text{\quad iff\quad} \Sat{M}{\Atom{\Obj A^2_0}{\Obj v_1, \Obj v_2}}[s].
\]
येथे चर-मूल्यांकनांचा समावेश करणे आवश्यक आहे हे लक्षात घ्या: आपण फक्त
``$Rab$ iff $\Sat{M}{\Atom{\Obj A^2_0}{a, b}}$'' असे म्हणू शकत नाही, कारण
$a$ आणि~$b$ ही आपल्या भाषेतील चिन्हे नाहीत; ती~$\Domain{M}$ ची
!!{element}s आहेत.

आपल्याकडे केवळ अणुसूत्र !!{formula}s नसून, तार्किक संयोजके आणि संख्यापक
वापरून ती एकत्र करता येतात; त्यामुळे अधिक गुंतागुंतीची !!{formula}s अशी इतर
संबंध परिभाषित करू शकतात जी~$\Struct M$ मध्ये थेट अंतर्भूत नसतात. ते कसे
करायचे आणि विशेषतः !!a{structure} मध्ये कोणते संबंध परिभाषित करता येतात,
यात आपल्याला रस आहे.
\end{explain}

\begin{defn}
$!A(\Obj v_1,\dots, \Obj v_n)$ हे~$\Lang L$ चे असे !!a{formula} असू द्या
की ज्यात फक्त~$\Obj v_1$,\dots, $\Obj v_n$ हीच चरे मुक्त येतात; आणि
$\Struct M$ ही~$\Lang L$ साठीची !!a{structure} असू द्या. कोणत्याही अशा
चर-मूल्यांकन~$s$ साठी की $s(\Obj v_i) = a_i$ ($i = 1, \dots, n$),
\[
Ra_1\dots a_n \text{\quad iff\quad} \Sat{M}{\Atom{!A}{\Obj
    v_1,\dots, \Obj v_n}}[s]
\]
असे असेल, तर आणि केवळ तेव्हाच $!A(\Obj v_1,\dots, \Obj v_n)$ हे
$R \subseteq \Domain M^n$ हा \emph{संबंध व्यक्त करते}.
\end{defn}

\begin{ex}
अंकगणिताच्या मानक प्रतिमान~$\Struct N$ मध्ये !!{formula} $\Obj
v_1 < \Obj v_2 \lor \eq[\Obj v_1][\Obj v_2]$ हे~$\Nat$ वरील~$\le$
संबंध व्यक्त करते. !!{formula} $\eq[\Obj v_2][\Obj v_1']$ हे उत्तरवर्ती
संबंध व्यक्त करते; म्हणजे, $m$ हा~$n$ चा उत्तरवर्ती असेल तेव्हा
$Rnm$ असणारा संबंध~$R \subseteq \Nat^2$. सूत्र
$\eq[\Obj v_1][\Obj v_2']$ हे पूर्ववर्ती संबंध व्यक्त करते. !!{formula}s
$\lexists[\Obj v_3][(\eq/[\Obj v_3][\Obj 0] \land
  \eq[\Obj v_2][(\Obj v_1 + \Obj v_3)])]$ आणि $\lexists[\Obj
  v_3][\eq[(\Obj v_1 + {\Obj v_3}')][\Obj v_2]]$ ही दोन्ही~$\Obj <$
संबंध व्यक्त करतात.
%READERNOTE{OLFOL-024: गोठवलेल्या इंग्रजीतील दुसऱ्या सूत्राच्या शेवटच्या v_2 पुढे वस्तुभाषेचे निर्देशचिन्ह चुकले आहे. समान परिच्छेदातील सर्व इतर चर-आस्थिती आणि अभिप्रेत सूत्राशी जुळण्यासाठी मराठीत वस्तुभाषेतील v_2 केले असून गोठवलेले इंग्रजी बाइट्स बदललेले नाहीत.}
याचा अर्थ अंकगणिताच्या भाषेत विधेयचिन्ह~$<$ प्रत्यक्षात अनावश्यक आहे; ते
परिभाषित करता येते.
\end{ex}

\begin{explain}
ही कल्पना फक्त विशिष्ट !!{structure}s मध्येच रंजक नाही; एखादे अभिप्रेत
प्रतिमान किंवा प्रतिमाने वर्णन करण्यासाठी आपण भाषा वापरतो तेव्हा, म्हणजे
उपपत्तींचा विचार करताना, ती सर्वसाधारणपणेही रंजक ठरते. अशा उपपत्तींमध्ये
मूलभूत चिन्हे म्हणून अनेकदा फक्त काही !!{predicate}s असतात, पण ती ज्या
प्रांताचे वर्णन करतात त्यात इतर अनेक संबंध महत्त्वाची भूमिका बजावतात. या
इतर संबंधांना भाषेतील मूलभूत !!{predicate}s चे अर्थनिर्धारण करणाऱ्या
संबंधांनी सुव्यवस्थित रीतीने व्यक्त करता येत असेल, तर आपण त्यांना त्या
भाषेत \emph{परिभाषित} करू शकतो असे म्हणतो.
\end{explain}

\begin{prob}
पुढील संबंध परिभाषित करणारी~$\Lang L_A$ मधील !!{formula}s शोधा:
\begin{enumerate}
\item $n$ हे~$i$ आणि~$j$ यांच्या दरम्यान आहे;
\item $n$ ने~$m$ ला निःशेष भाग जातो (म्हणजे $m$ ही~$n$ ची पटीत संख्या आहे);
\item $n$ ही अविभाज्य संख्या आहे (म्हणजे $1$ आणि~$n$ यांखेरीज कोणत्याही
  संख्येने~$n$ ला निःशेष भाग जात नाही).
\end{enumerate}
\end{prob}

\begin{prob}
सूत्र $!A(\Obj v_1, \Obj v_2)$ हे !!a{structure}~$\Struct M$ मध्ये संबंध
$R \subseteq \Domain M^2$ व्यक्त करते असे समजा. पुढील संबंध व्यक्त करणारी
सूत्रे शोधा:
\begin{enumerate}
\item $R$ चा व्युत्क्रम~$R^{-1}$;
\item सापेक्ष गुणाकार~$R \mid R$;
\end{enumerate}
$R$ चे संक्रमक संवरण~$R^+$ व्यक्त करण्याचा मार्ग तुम्हाला शोधता येतो का?
\end{prob}

\begin{prob}
फक्त एक द्विस्थानी विधेयचिन्ह~$<$ असलेली भाषा~$\Lang{L}$ असू द्या (इतर
कोणतीही !!{constant}s, !!{function}s किंवा !!{predicate}s नाहीत---अर्थात
~$\eq$ सोडून). $\Domain{N} = \Nat$ आणि
$\Assign{<}{N} = \Setabs{\tuple{n,m}}{n < m}$ असणारी रचना~$\Struct{N}$
असू द्या. पुढील विधाने सिद्ध करा:
\begin{enumerate}
\item $\{ 0 \}$ हा~$\Struct{N}$ मध्ये परिभाष्य आहे;
\item $\{ 1 \}$ हा~$\Struct{N}$ मध्ये परिभाष्य आहे;
\item $\{ 2 \}$ हा~$\Struct{N}$ मध्ये परिभाष्य आहे;
\item प्रत्येक $n \in \Nat$ साठी संच~$\{ n \}$ हा~$\Struct{N}$ मध्ये
  परिभाष्य आहे;
\item $\Domain{N}$ चा प्रत्येक सांत उपसंच~$\Struct{N}$ मध्ये परिभाष्य आहे;
\item $\Domain{N}$ चा प्रत्येक सांत-पूरक उपसंच~$\Struct{N}$ मध्ये परिभाष्य
  आहे (जिथे $X \subseteq \Nat$ हा सांत-पूरक असतो तेव्हा आणि केवळ तेव्हाच,
  जेव्हा $\Nat \setminus X$ हा सांत असतो).
\end{enumerate}
\end{prob}

\end{document}
